\chapter{Grundlagen}
\label{ch:grundlagen}

Dieses Kapitel beschreibt die fachlichen und systemtechnischen Grundlagen der
Arbeit. Zuerst wird das MXCK als Zielplattform eingeordnet. Danach werden die
für die Arbeit relevanten ROS~2-Begriffe erklärt. Abschließend wird der
Follow-the-Gap-Algorithmus sowohl in seiner ursprünglichen Form als auch in der
für dieses Projekt relevanten scan-basierten Variante erläutert.

\section{Das MXCarkit als Zielplattform}

Das MXCarkit ist ein Jetson-basiertes, containerisiertes Versuchsfahrzeug für
Lehre und Entwicklung im Bereich autonomer Fahrfunktionen. Laut der aktuellen
MXCK-Dokumentation werden die Systembestandteile bewusst in Docker-Container
aufgeteilt, wobei \texttt{mxck2\_ws} die Kernfunktionalität und
\texttt{development\_ws} experimentelle bzw.\ GPU-nahe Entwicklung aufnimmt.
Die zentrale Startlogik für Sensorik und Aktorik liegt im Paket
\texttt{mxck\_run}; einzelne Teilfunktionen wie LiDAR oder Kamera werden über
Launch-Argumente wie \texttt{run\_lidar:=true} aktiviert.

\begin{table}[H]
\centering
\caption{Relevante MXCK-Systemeigenschaften für diese Arbeit}
\label{tab:mxck_overview}
\begin{tabular}{>{\raggedright\arraybackslash}p{4.2cm} >{\raggedright\arraybackslash}p{9.2cm}}
\toprule
\textbf{Aspekt} & \textbf{Bedeutung für den FTG-Stack} \\
\midrule
NVIDIA Jetson & Rechnerplattform für ROS~2, Paket-Build, Sensorverarbeitung und Debugging \\
Docker-basierter Betrieb & Trennung von Sensor-/Control-Containern und Entwicklungsumgebung \\
2D-LiDAR (\texttt{/scan}) & Primärer Sensor für den reaktiven FTG-Pfad \\
Ackermann-Steuerung & Lenkwinkel- und Geschwindigkeitsvorgaben werden als Ackermann-Befehl ausgegeben \\
Bestehende Fahrzeugsteuerung & Der FTG-Stack darf die vorhandene VESC-/Safety-Kette nicht umgehen \\
\bottomrule
\end{tabular}
\end{table}

Für diese Arbeit ist vor allem wichtig, dass der autonome Output nicht direkt an
VESC-Topics gesendet wird, sondern an
\texttt{/autonomous/ackermann\_cmd}. Dadurch bleibt die bestehende MXCK-
Sicherheits- und Moduslogik erhalten.

\section{ROS~2: Workspace, Paket, Node und Launch}

Da der gesamte Stack auf ROS~2 basiert, ist ein sauberes Verständnis der
Architekturelemente notwendig.

\subsection{Workspace}

Ein ROS~2-Workspace ist ein Verzeichnis, in dem mehrere Pakete gemeinsam
gebaut werden. Gebräuchlich ist die Struktur
\texttt{src/}, \texttt{build/}, \texttt{install/}, \texttt{log/}. Im MXCK-Projekt
wird dieser Aufbau direkt genutzt: \texttt{mxck2\_ws} enthält die
Kernfunktionalität, während \texttt{development\_ws} für zusätzliche
Entwicklungsarbeit vorgesehen ist.

\subsection{Paket}

Ein Paket bündelt eine klar abgegrenzte Funktion. In dieser Arbeit werden
Pakete für Perception, Planner, Control und Bringup verwendet. Für Python-Pakete
sind insbesondere \texttt{package.xml}, \texttt{setup.py},
\texttt{setup.cfg}, der Python-Modulordner und optionale Launch-/Config-Ordner
relevant. Für C++-Pakete kommen stattdessen \texttt{CMakeLists.txt},
\texttt{package.xml}, Header- und Source-Dateien zum Einsatz.

\subsection{Node}

Ein Node ist ein laufender Prozess innerhalb des ROS~2-Graphs. Nodes abonnieren
Topics, veröffentlichen Topics, verwenden Parameter und werden meist über
Launch-Dateien gemeinsam gestartet.

\subsection{Launch-Dateien und YAML-Parameter}

ROS~2 nutzt bevorzugt Python-basierte Launch-Dateien. Diese starten mehrere
Nodes, reichen Parameter weiter und erlauben deklarative Startkonfigurationen.
Für dieses Projekt ist das besonders wichtig, weil die Pipeline nur dann stabil
und reproduzierbar ist, wenn Node-Namen, Executables, YAML-Dateien und Topics
konsistent zusammenpassen.

\section{Follow-the-Gap als reaktiver Planungsansatz}

Der Ursprung des Follow-the-Gap-Verfahrens liegt in der Arbeit von Sezer und
Gokasan. Dort wird FTG als lokaler Hindernisvermeidungsansatz formuliert, der
die Umgebung nur anhand der aktuell sichtbaren Hindernisse bewertet. Statt eine
globale Karte aufzubauen, sucht das Verfahren in jedem Zeitpunkt nach der
größten passierbaren Lücke und richtet das Fahrzeug möglichst auf deren Mitte
aus.

\subsection{Zentrale Idee der Originalarbeit}

Die Originalarbeit beschreibt drei Kernschritte:

\begin{enumerate}[label=\arabic*)]
    \item Aufbau einer Hindernis- bzw.\ Gap-Repräsentation,
    \item Bestimmung eines sicheren Gap-Center-Winkels,
    \item Fusion von Gap-Richtung und Zielrichtung zu einem finalen Heading.
\end{enumerate}

Wesentliche Eigenschaften des Verfahrens sind laut Originalarbeit:
\begin{itemize}
    \item Berücksichtigung des Sichtfelds,
    \item Berücksichtigung nicht-holonomer Fahrdynamik,
    \item keine lokale-Minimum-Problematik wie bei Potentialfeldern,
    \item nur wenige Tuning-Parameter.
\end{itemize}

\subsection{Gap-Idee in vereinfachter Form}

Im klassischen FTG werden aus den Hindernisgrenzen aufeinanderfolgende
``Gaps'' bestimmt. Die größte geeignete Lücke wird ausgewählt. Aus dieser
Lücke wird ein Winkel $\varphi_{gap}$ bestimmt. Anschließend wird dieser mit
einem Zielwinkel $\varphi_{goal}$ fusioniert, wobei nahe Hindernisse den
Einfluss des Gap-Winkels erhöhen und freie Umgebung den Einfluss des
Zielwinkels stärkt.

Formal beschreibt die Originalarbeit die finale Heading-Bildung als gewichtete
Fusion aus Gap-Winkel und Zielwinkel. Entscheidend ist dabei der kleinste
Hindernisabstand $d_{min}$: je kleiner $d_{min}$ ist, desto stärker soll die
Planung sicherheitsorientiert in Richtung des Gaps reagieren.

\subsection{Von der Originalformulierung zur praktischen Scan-Variante}

Die CTU-F1/10-Implementierung übernimmt die FTG-Idee in einen ROS-Node
\texttt{follow\_the\_gap}. Dessen klassischer ROS-Eingang ist ein
\texttt{sensor\_msgs/LaserScan} auf \texttt{/scan}; als Zielwinkel kann
optional \texttt{/lsr/angle} vorgegeben werden. Als wesentliche Ausgänge nennt
das CTU-Repository \texttt{/final\_heading\_angle},
\texttt{/visualize\_final\_heading\_angle},
\texttt{/visualize\_obstacles} und
\texttt{/visualize\_largest\_gap}.

Für das MXCK-Projekt wird dieser Kerngedanke nicht als
\texttt{/scan}-Direkteingang verwendet, sondern als \emph{scan-basierter
Vorverarbeitungspfad}:

\begin{lstlisting}
 /scan
   -> scan_preprocessor_node
   -> /autonomous/ftg/scan_filtered
   -> follow_the_gap_v0 (input_mode=scan)
   -> /final_heading_angle
   -> /gap_found
\end{lstlisting}

Damit wird der eigentliche C++-Planerkern entlastet: Das Re-Centering auf die
Fahrzeugfront, die Clipping-Logik und die Front-Clearance-Bestimmung finden
bereits vorher statt.

\section{Warum ROS~2 für dieses Thema relevant ist}

Die CTU-Arbeit von Endler zeigt, warum ROS~2 im Kontext schneller
autonomer Fahrfunktionen relevant ist. Zentrale Punkte sind:
\begin{itemize}
    \item DDS-basierte Kommunikation statt ROS~1-Master-Modell,
    \item klare Trennung von \texttt{rcl}, \texttt{rmw} und DDS,
    \item Build-System \texttt{ament} mit \texttt{colcon},
    \item moderne Python-Launchfiles,
    \item bessere Grundlage für Messung von Latenzen, Jitter und
    End-to-End-Ketten.
\end{itemize}

Für diese Projektarbeit ist das vor allem deshalb wichtig, weil das MXCK nicht
nur einen Algorithmus ``irgendwie ausführt'', sondern als Prozesskette aus
mehreren ROS~2-Nodes verstanden werden muss. Schon kleine Inkonsistenzen in
Topic-Namen, QoS, Frames oder Startreihenfolge können das Gesamtsystem
unbrauchbar machen.

\section{Zusammenfassung}

Die theoretische Grundlage dieser Arbeit ist damit zweigeteilt:
Einerseits liefert die FTG-Originalliteratur die mathematische und
planungslogische Basis. Andererseits liefert der aktuelle ROS~2-/F1TENTH-
Kontext die technische Form, in der diese Idee heute auf einem Jetson-basierten
Lehrfahrzeug umgesetzt wird. Genau an dieser Schnittstelle setzt die
vorliegende Arbeit an.
